\section{Method}

\begin{figure}[!htbp]
\centering
\includegraphics[width=\textwidth]{figures/method_framework.png}
\caption{Method framework for post-alignment NVD-GHSA field discrepancy typing and evidence-constrained adjudication. The adjudication stage is currently a silver-label prototype, not a human-gold source-truth result.}
\label{fig:method-framework}
\end{figure}


The method is an implemented baseline pipeline for post-alignment vulnerability-record fusion. Its input is an aligned pair \texttt{(r\_NVD, r\_GHSA)} keyed by CVE ID, and its unit of analysis is a target field within that pair. The pipeline has three operational stages: field-view construction, deterministic discrepancy typing, and an evidence-constrained source-support scaffold for sampled factual conflicts. The first two stages operate on all 8,066 aligned NVD-GHSA pairs. The third stage is evaluated with non-human annotations: silver-label development samples for severity and affected\_versions, plus two CVE-disjoint, prediction-sealed affected\_versions holdouts reviewed independently by Codex agents. It remains a controlled adjudication scaffold rather than a final source-of-truth system.

The design follows a type-first principle: triage before truth. The first decision asks what kind of disagreement is present. The second decision, activated only for residual factual conflicts, asks which source value has matching fetched-evidence cues, whether both values are provisionally defensible, or whether the case should be abstained for manual review. This separation is essential for this corpus because many binary differences are not contradictions: timestamp formatting, one-sided reference lists, and schema-level version-range differences require different downstream actions from incompatible severity labels.

The innovation claim is therefore deliberately narrow. The contribution is the post-CVE-alignment structured-field discrepancy task, the field-aware taxonomy and auditable baseline implementation, and the abstention-capable evidence scaffold for factual-conflict samples. The core method idea is the routing contract itself: a structured-field mismatch is first assigned an operational discrepancy type, and source-support analysis is delayed until the type indicates a residual factual conflict. It is not a new generic truth-discovery algorithm, not a final source-of-truth resolver, and not a semantic affected\_versions adjudicator.

The routing contract is summarized below. Its purpose is to make the routing decision explicit: the discrepancy type determines the next action before any source-support question is asked. This is the main difference from a winner-first fusion view, where every mismatch is treated as a conflict to be resolved.

\begin{table}[!htbp]
\centering
\setlength{\tabcolsep}{3pt}
\footnotesize
\caption{Method contract for routing field discrepancies before source-support adjudication. Validation boundaries reflect the current validation state.}
\label{tab:03-method-table-1}
\begin{tabularx}{\textwidth}{>{\raggedright\arraybackslash}p{0.18\textwidth}>{\raggedright\arraybackslash}X>{\raggedright\arraybackslash}X>{\raggedright\arraybackslash}X}
\toprule
Discrepancy type & Typical trigger in this study & Downstream action & Current validation boundary \\
\midrule
equivalent & Canonical severity, CWE set, or affected-version span matches & Merge or suppress as non-actionable difference & Deterministic baseline; pending RQ2 gold check \\
representation & Same-day timestamps, label aliases, URL/range-shape formatting & Normalize representation and keep provenance & Deterministic baseline; pending RQ2 gold check \\
incomplete & One-sided references, missing CWE IDs, one-sided affected-version content & Treat as augmentation or coverage gap, not source selection & Deterministic baseline; pending RQ2 gold check \\
temporal & Cross-day published/date differences & Interpret as publication or enrichment timing & Deterministic baseline; pending RQ2 gold check \\
factual conflict & Residual incompatible severity labels or unresolved version structures & Construct evidence context and allow NVD, GHSA, both, neither, or abstain & RQ3 prototype only; human audit pending \\
\bottomrule
\end{tabularx}
\end{table}

\subsection{Operational input/output contract}

The method operates on field instances rather than whole records. A field instance is \path{i = (cve_id, f, r_NVD.f, r_GHSA.f)}, where \texttt{f} is one target field and the two records have already been aligned by CVE ID. The implementation does not overwrite either source record. Each stage appends a new record around the original values so that later analysis can trace a decision back to source values, normalization, rule triggers, evidence URLs, and human-audit fields.

\begin{table}[!htbp]
\centering
\setlength{\tabcolsep}{3pt}
\footnotesize
\caption{Operational input/output contract for the implemented method stages. Validation boundaries reflect the current validation state.}
\label{tab:03-method-operational-input-output-contract-1}
\begin{tabularx}{\textwidth}{>{\raggedright\arraybackslash}p{0.18\textwidth}>{\raggedright\arraybackslash}X>{\raggedright\arraybackslash}X>{\raggedright\arraybackslash}X}
\toprule
Stage output & Required content & Used for & Current boundary \\
\midrule
Field view \texttt{V\_i} & Raw source values, normalized comparands, source provenance, rule-note carrier & RQ1 distribution and RQ2 typing audit & Normalization is baseline logic, not validated semantics \\
Typed instance \texttt{T\_i} & \texttt{V\_i}, discrepancy label, rule identifier or rule note, routing action & RQ1 counts, RQ2 rule-trigger diagnostic, FC sampling & Labels are deterministic outputs until RQ2 gold exists \\
Evidence context \texttt{E\_i} & Candidate URLs, URL provenance, fetch status, fetched snippets, source-class diagnostics & RQ3 evidence availability and audit worklist & Snippets are evidence inputs, not verified page reconstructions \\
Adjudication record \texttt{A\_i} & Source-support decision, matched cues, abstention state, required human-audit fields & RQ3 silver diagnostic and later human-gold evaluation & Silver labels are annotation aids, not gold labels \\
Sealed holdout record \texttt{H\_i} & CVE-disjoint cohort commitment, blind evidence worklist, prediction-only seal, dual-review decision, strict consensus & RQ3 affected\_versions leakage and task-definition audit & Dual-Codex consensus is non-human and selective; human signoff remains pending \\
\bottomrule
\end{tabularx}
\end{table}

Procedure 1 runs over the aligned corpus and is deterministic. For each CVE-aligned pair and target field, it preserves both source values, builds field-specific normalized comparands, applies the ordered rules described below, and emits a typed field instance with a rule note and routing action. Non-FC labels stop at the routing layer: they are not passed to source-support adjudication. Procedure 2 is activated only for sampled factual-conflict candidates. It collects candidate URLs from NVD/GHSA reference contexts, records URL provenance and fetch status, extracts text snippets when available, matches diagnostic field-specific cues, and emits a provisional NVD, GHSA, both, neither, or abstain decision for human audit. The guarded evaluator computes human-gold metrics only after valid final human rows exist.

This contract is also the reproducibility boundary. The paper-facing artifacts can show which field instance was processed, which rule routed it, which evidence was available, and which final human fields remain blank. A claim is allowed only if the corresponding stage output exists and the relevant validation boundary has been satisfied.

\subsection{Field-view construction}

The field-view stage extracts the target fields from each aligned pair and normalizes them into comparable representations. The implemented baseline covers severity, published/date, references, affected\_versions, and cwe\_ids. Each field view preserves the original source values, the normalized comparison representation, and a rule note explaining the path taken by the deterministic comparator. This makes the later diagnostics auditable: a field instance can be traced back to the source records, the normalization step, and the exact rule trigger used by the baseline.

Severity labels are mapped to canonical CVSS-style labels so that source-specific labels such as moderate and medium compare under the same vocabulary. Published/date values are parsed and compared after timestamp normalization, while preserving the distinction between same-day representation differences and cross-day timing differences. References are deduplicated by URL and host-level structure so that subset and formatting effects can be separated from residual conflicts. CWE identifiers are canonicalized as sets. Affected\_versions values are converted into span-like records when the source schema allows it, with the current implementation retaining range-bound shape, endpoint strings, and package/context information for later diagnostics.

For references, the development workflow treats URL normalization as a resource-identity hypothesis rather than a string-cleaning fact. A complete-impact audit groups only the URL identities changed by a candidate profile, freezes live-HTTP metadata and content fingerprints, masks the candidate transformation and expected label from two reviewers, and retains a change only under strict two-run agreement. The audited development profile keeps exact HTTP redirects, scoped Liferay presentation-query removal, and identifier-preserving GitHub-advisory and Huntr aliases, but excludes encoded GitHub line-suffix stripping after the reviewers apply incompatible underlying-file and HTTP-resource interpretations. This profile is selected after the audit and is therefore a post-hoc development artifact, not a confirmed normalization rule.

The field-view stage is intentionally conservative. It does not try to infer missing package identity or repair upstream record alignment. As a diagnostic baseline, when affected\_versions values use incompatible package names, ecosystem conventions, or range schemas, the current implementation records those structures for audit rather than silently forcing them into a shared range model.

Validation status: field-view construction is implemented and reproducible over the aligned corpus. Its normalization choices are still baseline rules and should be evaluated through the pending RQ2 human labels before being treated as validated field semantics.

\subsection{Discrepancy typing}

The discrepancy-typing stage applies field-specific deterministic rules and emits one of the five working labels: equivalent, representation\_discrepancy, incomplete, temporal\_discrepancy, or factual\_conflict. Equivalent means the normalized values carry the same field meaning. Representation\_discrepancy means the raw strings or schemas differ but the implemented rule treats the intended field fact as compatible. Incomplete means one side has field content that the other side lacks. Temporal\_discrepancy captures publication/date differences that are better modeled as process timing than as a contradictory value. Factual\_conflict is the residual category used when the implemented rules do not find equivalence, representation compatibility, incompleteness, or temporal explanation.

The rules are field-specific because the same surface operation has different meaning across fields. The general priority is conservative: preserve missingness first when one side lacks comparable content; test field-specific equivalence and representation compatibility; route one-sided coverage gaps to incompleteness; route publication timing differences to temporal\_discrepancy; and use factual\_conflict only as the residual category. When multiple non-conflict explanations are possible, the implementation emits the most operationally specific routing label: temporal for date timing, incomplete for one-sided content, representation\_discrepancy for compatible form/schema differences, and equivalent for normalized identity.

\begin{table}[!htbp]
\centering
\setlength{\tabcolsep}{3pt}
\footnotesize
\caption{Ordered field-specific baseline rule paths for discrepancy typing. The rules are diagnostic until RQ2 human labels exist.}
\label{tab:03-method-discrepancy-typing-1}
\begin{tabularx}{\textwidth}{>{\raggedright\arraybackslash}p{0.18\textwidth}>{\raggedright\arraybackslash}X>{\raggedright\arraybackslash}X}
\toprule
Field & Ordered baseline rule path & Boundary or known limitation \\
\midrule
severity & Canonicalize CVSS-style labels; matching canonical labels route away from conflict; missing one-sided values route to incomplete; different canonical labels become factual-conflict candidates. & Label canonicalization is deterministic; correctness of a source rating is not decided here. \\
published/date & Parse timestamps; same-day differences route to representation\_discrepancy; cross-day differences route to temporal\_discrepancy. & Date routing explains process timing, not which publication timestamp is more authoritative. \\
references & Normalize URLs and host-level structure; URL-format variants route to representation\_discrepancy; strict one-sided or subset reference sets route to incomplete; rare residual disagreement remains factual\_conflict. & Reference provenance is not evidence correctness; a URL's source side does not decide truth. \\
cwe\_ids & Canonicalize CWE IDs as sets; set equality routes to equivalent; one-sided or subset content routes to incomplete; residual set disagreement remains factual\_conflict. & CWE granularity and upstream mapping policy are not adjudicated by this baseline. \\
affected\_versions & Preserve package/context, endpoint strings, range-bound shape, and span-like records; baseline structural or effective-bound heuristics can route cases away from conflict; one-sided content routes to incomplete; residual incompatible version structures remain factual-conflict candidates. & Current logic is structural and diagnostic; it is not package-aware or range-aware semantic adjudication. \\
\bottomrule
\end{tabularx}
\end{table}

Examples are used only to clarify diagnostic-only routing behavior. A severity pair such as \texttt{HIGH} versus \texttt{CRITICAL} remains a factual-conflict candidate; a same-day timestamp string difference is a representation discrepancy; cross-day publication values are temporal; a strict-subset reference list is incomplete; and an NVD point version inside a parseable GHSA range is routed as baseline range-shape compatibility, pending human check. For affected\_versions, this remains an audit-prioritization heuristic and diagnostic token-support prototype: package names that do not overlap and version tokens that do not align are kept as audit-risk residuals rather than treated as semantic version-range adjudication.

Affected\_versions requires the most structural handling. The implemented baseline compares span-like records and applies range-bound diagnostics such as effective-bound equivalence, one-sided presence, shared upper-bound endpoints, and the 2026-04-22 rule tightenings that lowered affected\_versions factual-conflict candidates from 1,272 to 652. The 2026-07-15 input-integrity refresh then filtered NVD CPE matches explicitly marked \texttt{vulnerable=false}, lowering the current residual count from 652 to 651. These changes are implementation diagnostics: they show how many cases the pipeline routes differently, not that correctness improved by the same amount. The tightenings include cases where \texttt{end\_including} and \path{end_excluding/fixed} share a major.minor prefix, where an NVD point version falls inside a parseable GHSA range, where \texttt{end\_excluding} strings differ only by prefix truncation, and where NVD uses a date string as a version-end boundary. These rules are audit-prioritization heuristics, not human-validated truth. They can fail when the two sources name different packages, use ecosystem-specific ordering, encode branch- or distribution-specific backports, or express the same advisory through incompatible range schemas. Because the planned 300-instance discrepancy-typing gold set is not complete, the current RQ2 artifact reports rule-trigger distributions and audit targets rather than typing accuracy, precision, recall, or macro-F1.

The diagnostic input-integrity analysis is separately recorded at \path{results/input_integrity/affected_versions/affected_versions_input_integrity.{json,md\}}. The legacy aligned corpus contained 1,105 \texttt{vulnerable=false} CPE items across 106 of the 8,066 matched records. Removing them changed 10 affected\_versions baseline classifications. Seven rows in the frozen 100-sample affected\_versions cohort had refreshed source inputs, but the diagnostic cohort mapping, its 40 determinate versus 60 abstained AI-adjudicated decisions, and all final labels remained unchanged after targeted recheck. This recheck is AI provenance (\texttt{label\_is\_human=false}), not human validation. A scan of the current 28,785-record reviewed GHSA snapshot found no range containing multiple events, so the present first-event flattening did not alter this snapshot; it remains a latent risk for future snapshots rather than a validated general assumption.

Validation status: discrepancy typing is implemented and useful for audit prioritization. It is not yet a gold-backed classifier, so all RQ2 accuracy and agreement claims remain pending.

\subsection{Evidence context construction}

The evidence stage is activated only for sampled factual-conflict candidates. It collects candidate URLs from the aligned records' reference contexts, fetches text snippets and basic metadata, and stores a per-sample evidence context. The evidence context keeps source provenance separate from source truth: a URL can originate from NVD references only, GHSA references only, or both, but that provenance alone does not decide which field value is correct.

For the severity seed, the project assembled 80 examples, collected 470 candidate evidence URLs, and cached fetched evidence text. Of those URL records, 415 returned HTTP-ok fetched text snippets, 46 hit URL errors, 4 returned 403, 3 returned 404, and 2 timed out. All 80 samples have at least one successful text evidence source. For affected\_versions, the same evidence builder has been run on the 100-sample manual-check set. Its manifest records 585 candidate evidence URLs: 402 HTTP-ok responses, of which 401 yielded OK text records after one skipped non-text URL, plus 132 URL errors, 30 timeouts, 15 HTTP 403 responses, 3 HTTP 404 responses, and 2 HTTP 308 responses. All 100 samples have at least one successful text evidence source.

The evidence-source reliability diagnostic further classifies each fetched URL by source class and reference provenance. Source classes distinguish NVD pages, GitHub advisories, GitHub commits or repositories, vendor advisories, package registries, mailing lists or issue trackers, and other pages. These classes are heuristic audit features, not source-truth labels. They are used to characterize evidence availability, identify fetch bias, and prioritize manual review when a field's evidence comes mainly from pages that are hard to fetch or weakly structured.

Validation status: evidence-context construction is implemented for the current severity and affected\_versions samples. The fetched snippets are evidence inputs, not manually verified page reconstructions, and fetchability differences remain a validity risk.

\subsection{Evidence-constrained source-support prototype}

The adjudication decision space is source-support oriented rather than winner-only. For human-audited evaluation, a sampled factual conflict can be labeled as NVD-supported, GHSA-supported, both-source supported, neither-source supported, or abstain. NVD-supported and GHSA-supported mean that the available evidence is judged to support one side's field value. Both-source supported means that both values are defensible or supported under different source contexts. Neither-source supported means that the fetched evidence does not support either recorded value. Abstain means that the evidence is missing, weak, contradictory, fetch-biased, or outside the prototype's cue model. The prototype output therefore records the predicted source decision, matched evidence cues, and an abstention path rather than forcing a value in every case.

The current evidence-aware silver-v2 label sets are generated from the fetched evidence contexts and record provisional source decisions for each sample, including both-source support, NVD-only, GHSA-only, neither-source, and abstain outcomes. Silver-v2 is an annotation aid and diagnostic comparison target, not an evaluation oracle. The severity prototype uses an \texttt{evidence\_score\_baseline} that matches severity labels, numeric scores, and CVSS vectors in fetched text snippets. The affected\_versions prototype uses a \path{version_token_support_baseline} that checks whether fetched text contains version tokens from either side. The latter is intentionally shallow: it does not parse package identity, ecosystem-specific version ordering, or version-range containment, which are known schema-sensitive parts of OSV/GHSA affected-version representation and affected-version identification \cite{OSVSchema2026,GHSARepo2026,Chen2025AffectedVersions}.

Fixed-source baselines (\texttt{prefer\_nvd}, \texttt{prefer\_ghsa}) and a \texttt{latest\_published} baseline provide simple comparisons. Within the silver-label setup, the evidence-aware baselines test whether lightweight evidence matching changes the diagnostic pattern relative to fixed-source preferences, while preserving the distinction between supported, unsupported, and abstained cases. The current results should therefore be read as selective baselines against silver labels, not as final adjudication performance.

For affected\_versions, the v1 holdout protocol excludes every CVE in the development cohort before fixed-hash sampling, uses a dedicated evidence cache, strips prior labels and predictions from the reviewer worklist, and seals label-free method predictions before review. Two Codex reviewers independently assign both a discrepancy type and source support. The merger accepts only determinate exact joint agreement that satisfies the evidence-support contract. Its preregistered endpoint evaluates source selection on all accepted rows; a post-unseal diagnostic then separates true factual conflicts from non-conflicts after discovering that the original endpoint conflates the two tasks.

The v2 protocol turns that failure into an explicit design constraint. It excludes both the development and v1 CVEs, preregisters discrepancy typing and source adjudication on strict factual conflicts as separate endpoints, and evaluates an independent source head without gating it on the method's own type prediction. Reviewer evidence citations are bound to literal frozen snippets through structured endpoint, target, and role claims. Prediction files, method code, protocol code, reviewer identities, row order, and evidence hashes are sealed or checked before evaluation. This protocol prevents the v1 metric conflation; it does not by itself validate either predictor.

Validation status: the adjudication stage is a prototype. Its current role before human audit is to define the source-support decision space, exercise evidence availability, and generate a worklist for manual adjudication.

\subsection{Audit interface and current boundary}

The implementation already prepares the handoff from prototype adjudication to human audit. The RQ2 annotation seed and 20\% consistency-review template cover the target fields and rule-trigger strata, but their labels remain blank. The RQ3 gold-audit templates preserve source identifiers, fetched evidence, silver-v2 decisions, and required human-audit fields for 80 severity rows and 100 affected\_versions development rows, but all final human audit rows are still draft. The guarded RQ3 evaluator refuses to write gold-audit metrics until valid final human rows exist. The affected\_versions holdout consensus rows are also marked \texttt{label\_is\_human=false} and cannot substitute for that audit.

The method stages map to the research questions with different maturity levels. RQ1 uses Procedure 1 over all aligned pairs and target fields to report a deterministic discrepancy distribution. RQ2 uses the same typed instances plus annotation templates and supports rule-trigger/readiness diagnostics together with bounded non-human full-impact audits for references and CWE candidates. RQ3 uses Procedure 2 over sampled factual-conflict candidates and currently supports evidence availability, non-human label diagnostics, holdout agreement/coverage, and evaluation-protocol findings rather than human-gold adjudication performance.

This boundary is part of the method rather than a footnote. The current paper can describe the field-view construction, deterministic typing rules, bounded non-human full-impact audits, evidence-context builder, sealed holdout protocol, abstention mechanism, and audit interface. It cannot yet claim human-gold RQ2 typing performance, validated affected\_versions false-positive rates, an affected\_versions method improvement, or final RQ3 source-truth adjudication performance.
